\chapter{Minkowski identities and Euclidean CMC rigidity}\label{sec:cmc-minkowski}
The volume constraint admits a particularly useful test function built from the support function. We derive the integral identities and then identify an exactly volume-preserving path producing that test. The rigidity theorem is due to Barbosa and do Carmo; \cite{Wen91}, pp.~375--379, gives an elementary parallel-surface proof. Our convention is $A=D\nu$ and $H=\operatorname{tr}A$, so an outward sphere has $H=n/R$. Wente's scalar mean curvature is $-H/n$ in this notation.

\begin{definition}[Support function and tangential position]\label{def:cmc-support-function}
For a smooth oriented hypersurface $\Sigma^n\subset\mathbb R^{n+1}$, with position vector $x$ relative to a fixed origin and unit normal $\nu$, its support function and tangential position field are
\[
 u=\langle x,\nu\rangle,\qquad x^T=x-u\nu.
\]
Neither $u$ nor $x^T$ requires star-shapedness; $u$ may change sign. The variation used in \cite{Wen91}, pp.~377--378, combines normal displacement with dilation and thus involves $u$.
\end{definition}

\begin{lemma}[Tangential position and the first Minkowski identity]\label{lem:cmc-first-minkowski}
For the notation of Definition \ref{def:cmc-support-function},
\[
 \nabla_\Sigma x^T=I-uA,\qquad \operatorname{div}_\Sigma x^T=n-Hu.
\]
If $\Sigma$ is compact without boundary, then $\int_\Sigma Hu\,dA=n\operatorname{Area}(\Sigma)$. This is the integral identity underlying \cite{Wen91}, equation (10) and its appendix.
\end{lemma}
\begin{proof}
For a tangent vector $v$, differentiate $x-u\nu$ in Euclidean space and project tangentially. The derivative of $x$ is $v$, the term $v(u)\nu$ has zero tangential projection, and the remaining term is $-uAv$. Taking the trace gives the divergence formula. The divergence theorem on the closed hypersurface gives the integral identity. No sign condition on $u$ or on the principal curvatures enters this argument.
\end{proof}

\begin{lemma}[The first Newton tensor has zero divergence]\label{lem:cmc-newton-tensor}
On a Euclidean hypersurface put $P_1=HI-A$. Then
\[
 \operatorname{div}_\Sigma P_1=0,\qquad
 \operatorname{tr}P_1=(n-1)H,\qquad
 \operatorname{tr}(P_1A)=H^2-|A|^2.
\]
Here $H$ need not be constant.
\end{lemma}
\begin{proof}
The Euclidean Codazzi equation $\nabla_i h_{jk}=\nabla_j h_{ik}$ follows by commuting ambient derivatives of the normal and using zero ambient curvature. Contracting the first and third indices gives $\nabla^i h_{ij}=\nabla_jH$. Thus the divergences of $HI$ and $A$ cancel. The trace formulas follow by expanding $P_1$ in an orthonormal principal frame. These are local differential identities, independent of stability or compactness.
\end{proof}

\begin{proposition}[The second Minkowski identity]\label{prop:cmc-second-minkowski}
Every smooth closed oriented Euclidean hypersurface satisfies
\[
 \int_\Sigma u(H^2-|A|^2)\,dA=(n-1)\int_\Sigma H\,dA.
\]
In particular, if $H$ is constant and $a=\operatorname{Area}(\Sigma)$, then
\[
 \int_\Sigma u|A|^2\,dA=Ha.
\]
This supplies a divergence-form route to the curvature defect in \cite{Wen91}, equation (11).
\end{proposition}
\begin{proof}
Use Lemma \ref{lem:cmc-newton-tensor} and Lemma \ref{lem:cmc-first-minkowski} to compute
\[
 \operatorname{div}_\Sigma(P_1x^T)
 =\operatorname{tr}(P_1\nabla x^T)
 =(n-1)H-u(H^2-|A|^2).
\]
Its integral vanishes. When $H$ is constant, the first Minkowski identity gives $H\int u=na$. Therefore $H^2\int u-(n-1)Ha=Ha$, which is the claimed second formula. This deduction does not divide by $H$.
\end{proof}

\begin{proposition}[The support function solves an inhomogeneous Jacobi equation]\label{prop:cmc-support-jacobi}
For a Euclidean hypersurface, let $J=\Delta_\Sigma+|A|^2$. Then
\[
 Ju=H+\langle\nabla H,x^T\rangle.
\]
Thus $Ju=H$ when $H$ is constant. Compare the dilation component of the variation in \cite{Wen91}, p.~378.
\end{proposition}
\begin{proof}
Differentiating $u=\langle x,\nu\rangle$ gives $\nabla u=Ax^T$. Taking its divergence, contracting Codazzi as above, and using $\nabla x^T=I-uA$ yields
\[
 \Delta u=\langle\nabla H,x^T\rangle+\operatorname{tr}(A(I-uA))
 =\langle\nabla H,x^T\rangle+H-u|A|^2.
\]
Move the last term to the left. In the CMC case this also agrees with the normal-variation formula: under dilation the initial normal speed is $u$, while the derivative of $H$ is $-H$.
\end{proof}

\begin{proposition}[CMC area-volume balance]\label{prop:cmc-area-volume-balance}
Let $\Sigma=\partial\Omega\subset\mathbb R^{n+1}$ be a nonempty smooth compact boundary of a bounded region of positive volume, with outward normal and constant $H$. Then
\[
 H(n+1)\operatorname{Vol}(\Omega)=n\operatorname{Area}(\Sigma).
\]
Consequently $H>0$. This is \cite{Wen91}, equation (10), in the outward, unaveraged convention.
\end{proposition}
\begin{proof}
The ambient divergence theorem applied to $x$ gives $\int_\Sigma u=(n+1)\operatorname{Vol}(\Omega)$. Insert this into the first Minkowski identity. Positivity of volume and area gives the sign of $H$. More generally, the first Minkowski identity alone rules out any nonempty closed Euclidean minimal hypersurface, without needing an enclosed region.
\end{proof}

\begin{proposition}[A zero-mean test measures the full umbilicity defect]\label{prop:cmc-umbilicity-test}
Let $\Sigma$ be a smooth closed Euclidean CMC hypersurface. Define
\[
 f=1-\frac Hn u,\qquad D=|A|^2-\frac{H^2}{n}.
\]
Then $\int_\Sigma f=0$, $Jf=D$, and the Euclidean second-variation form satisfies
\[
 Q(f,f)=-\int_\Sigma D\,dA.
\]
The function $f$ is the normal speed identified in \cite{Wen91}, p.~378.
\end{proposition}
\begin{proof}
The first Minkowski identity gives the zero mean. Since $J1=|A|^2$ and $Ju=H$, the Jacobi equation follows. The first and second Minkowski identities also give
\[
 \int uD=Ha-\frac{H^2}{n}\int u=0.
\]
Integrating by parts in the second-variation formula of Proposition \ref{prop:cmc-second-variation}, we obtain
$Q(f,f)=-\int fJf=-\int D+(H/n)\int uD=-\int D$.
Pointwise Cauchy--Schwarz gives $D\geq0$, with equality exactly when $A=(H/n)I$. All integrals use the induced area measure.
\end{proof}

\begin{proposition}[Parallel displacement followed by volume correction]\label{prop:cmc-parallel-volume-path}
Under the hypotheses of Proposition \ref{prop:cmc-area-volume-balance}, the test in Proposition \ref{prop:cmc-umbilicity-test} is realized by an exactly volume-preserving family
\[
 y_t=s(t)(x+t\nu),\qquad
 s(t)=\left(\frac{\operatorname{Vol}(\Omega)}{\operatorname{Vol}(\Omega_t)}\right)^{1/(n+1)},
\]
where $\partial\Omega_t=\{x+t\nu:x\in\Sigma\}$ for sufficiently small $|t|$. This is the family in \cite{Wen91}, Section II, with signs translated to our convention.
\end{proposition}
\begin{proof}
The tubular neighborhood theorem for the compact embedded boundary makes $x+t\nu$ embedded for small $|t|$ and gives a smoothly varying bounded region $\Omega_t$. In a principal frame the differential is $I+tA$, so its area element is $\prod_i(1+t\kappa_i)dA$. For sufficiently small $t$ these factors are positive, and its outward normal remains $\nu$. The volume transport formula gives $V'(0)=a$, where $V(t)=\operatorname{Vol}(\Omega_t)$ and $a=\operatorname{Area}(\Sigma)$.
Dilation by the positive factor $s(t)$ multiplies volume by $s(t)^{n+1}$, so the corrected volume is exactly $V(0)$. Also
\[
 s'(0)=-\frac{a}{(n+1)V(0)}=-\frac Hn.
\]
Thus $y'_0=\nu-(H/n)x$ has normal part $f\nu$. Its tangential part is harmless for the second variation at a CMC critical point, as established in Proposition \ref{prop:cmc-second-variation}. Consequently the second derivative of area along this explicit path is $-\int D$.
\end{proof}

\begin{theorem}[Euclidean volume-stability rigidity]\label{thm:euclidean-cmc-stability-rigidity}
Let $n\geq2$ and let $\Sigma=\partial\Omega\subset\mathbb R^{n+1}$ be a smooth connected compact embedded CMC boundary. It is stable under the total-volume constraint if and only if it is a round sphere. This is the embedded-boundary case of the Barbosa--do Carmo theorem, stated and proved by a parallel-family argument in \cite{Wen91}, pp.~375--379.
\end{theorem}
\begin{proof}
If $\Sigma$ is volume-constrained stable, Proposition \ref{prop:cmc-umbilicity-test} gives $0\leq Q(f,f)=-\int D$. Since $D\geq0$ is continuous, $D=0$ everywhere. Hence $A=(H/n)I$. Proposition \ref{prop:cmc-area-volume-balance} gives $H>0$. For every tangent vector $v$,
\[
 D_v\left(x-\frac nH\nu\right)=v-\frac nH Av=0.
\]
Connectedness gives a fixed center $b=x-(n/H)\nu$. Thus $|x-b|=n/H$, and $\Sigma$ lies in the corresponding sphere. Its inclusion is locally open in that sphere because the tangent spaces have the same dimension; its compact image is closed. The sphere is connected, so the image is the whole sphere. Conversely Proposition \ref{prop:round-sphere-poincare} proves the required nonnegativity for every zero-mean speed on a round sphere. Proposition \ref{prop:realize-volume-constraint} identifies this condition with stability along exact volume-preserving variations. No star-shapedness or pointwise convexity is assumed in the argument.
\end{proof}
